% Clase 14 --- El experimento de Oersted (§2.2 y §2.4).
% El docente construyó el argumento sustituyendo una cosa por otra: primero un
% imán grandote con imancitos livianos alrededor, y después "en lugar de poner
% un imán grandote acá, lo que hizo fue poner un cable con una corriente".
% Los dos paneles tienen los mismos imancitos: lo único que cambia es qué los
% orienta.
\providecommand{\imancito}[3]{%
  \begin{scope}[shift={(#1,#2)}, rotate=#3]
    \fill[figblue!35] (-0.28,-0.09) rectangle (0,0.09);
    \fill[figred!35] (0,-0.09) rectangle (0.28,0.09);
    \draw[figgray, line width=0.5pt] (-0.28,-0.09) rectangle (0.28,0.09);
  \end{scope}}
\begin{center}
\begin{tikzpicture}[scale=1]

  % =============== (a) un imán grandote ===============
  \begin{scope}
    % el imán, visto desde arriba
    \fill[figblue!18] (-0.5,-1.3) rectangle (0.5,0);
    \fill[figred!18] (-0.5,0) rectangle (0.5,1.3);
    \draw[figgray, line width=1pt] (-0.5,-1.3) rectangle (0.5,1.3);
    \draw[figgray, line width=0.7pt] (-0.5,0) -- (0.5,0);
    \node[figlbl, figblue] at (0,-0.65) {S};
    \node[figlbl, figred] at (0,0.65) {N};

    % imancitos livianos alrededor, orientados por el grande
    \imancito{-1.7}{0.9}{-35}  \imancito{1.7}{0.9}{215}
    \imancito{-1.9}{0}{-90}    \imancito{1.9}{0}{-90}
    \imancito{-1.7}{-0.9}{35}  \imancito{1.7}{-0.9}{145}

    \node[figlbl, anchor=north, align=center] at (0,-2.1)
      {(a) un imán grande\\[-0.2em] \footnotesize (por ejemplo, la Tierra)};
  \end{scope}

  % =============== (b) un cable con corriente ===============
  \begin{scope}[xshift=7.6cm]
    % el cable, perpendicular a la mesa: se ve el corte
    \draw[figred, line width=1.1pt] (0,0) circle (0.17);
    \node[figcarga, fill=figred, minimum size=3.4pt] at (0,0) {};
    \node[figlbl, figred, anchor=north, scale=0.92] at (0,-0.4) {$I$ sale del plano};

    % los mismos imancitos, ahora tangentes a las circunferencias
    % tangentes, en sentido antihorario: con I saliendo del plano, el N de cada
    % imancito apunta según B, o sea a 90 grados de su posición angular.
    \imancito{-1.55}{0.9}{240}  \imancito{1.55}{0.9}{120}
    \imancito{-1.8}{0}{270}     \imancito{1.8}{0}{90}
    \imancito{-1.55}{-0.9}{300} \imancito{1.55}{-0.9}{60}
    \draw[figaux] (0,0) circle (1.8);

    \node[figlbl, anchor=north, align=center] at (0,-2.1)
      {(b) un cable con corriente\\[-0.2em] \footnotesize (Oersted)};
  \end{scope}

\end{tikzpicture}\\[0.5em]
{\small\itshape Los imancitos livianos hacen de detector: se orientan solos
según lo que haya alrededor, mientras que el objeto del centro está fijo. El
punto del experimento es que en (b) reaccionan igual que en (a), aunque ahí no
hay ningún imán: alcanza con que por el cable circule corriente ---con las
cargas quietas no pasa nada---. Y el efecto es recíproco: una espira con
corriente colocada cerca de un imán se orienta y es atraída como si fuera otro
imán, que es exactamente un electroimán. La conclusión que ordena todo el
capítulo es que \textbf{los efectos magnéticos los producen cargas en
movimiento}, y que lo que siente fuerza magnética son también cargas en
movimiento. Los imanes permanentes no son la excepción: a nivel molecular hay
corrientes diminutas que en casi todos los materiales apuntan al azar, y en un
imán están alineadas y sus efectos se suman. Notar que la orientación en (b) es
\emph{tangente} a las circunferencias alrededor del cable, no radial.}
\end{center}
